\documentclass{article}
\usepackage{amsmath}
\usepackage{hyperref}
\usepackage{amssymb}
\usepackage{dsfont}

\author{Loke Strøm Haugsnes}
\setlength{\parindent}{0pt}

\title{Mathematics for Drone Piloting} 

\begin{document}
\maketitle

This document describes the mathematical model used for rigid body dynamics
in the drone piloting simulation.

\section{Rigid body}

\subsection*{Definitions:}

\begin{tabbing}
	\hspace{1cm}\= \kill
	\> \( M \) is the mass of the rigid body. \\
    \> \( P \) is the position of the rigid body. \\
    \> \( v \) is the velocity of the rigid body. \\
    \> \( a \) is the acceleration of the rigid body. \\
	\> \( O_M \) is the center of mass of the rigid body in body frame. \\
	\> \( I \) is the inertia tensor of the rigid body about its center of mass. \\
	\> \( L \) is the angular momentum of the rigid body. \\
    \> \( q \) is the orientation of the rigid body, represented as a quaternion. \\
	\> \( \tau \) is the torque applied to the rigid body. \\
    \> \( F_S \) is the sum of the forces applied to the rigid body in body frame. \\
    \> \( F_{DS} \) is the sum of the force displacement from \( O_M \). \\
\end{tabbing}
The subscript \( B \) denotes that the variable is expressed in the body frame.

\subsection*{Variable Initialization:}
From the drone configuration file:
\begin{tabbing}
    \hspace{1cm}\= \kill
    \> \( M \) is defined as the mass of the drone. \\
    \> \( I_B \) is defined as the inertia tensor of the drone. \\
    \> \( O_M \) is defined as the center of mass of the drone. \\
\end{tabbing}
Note: Other variables in the JSON configuration file have no effect on the math, 
as they are only used for visualization and other purposes even the motor properties, 
as they are used to calculate the forces applied to the drone, which is represented by \( F_S \) and \( F_{DS} \). \\
Then the rest of the variables will be initialized too zero, as the drone starts at rest.
Or they may be affected by the map configuration file or other factors like swapping the drone mid flight.
This covers all the variables initialized state. 
The next step is to define how these variables are updated during the simulation. 
But before that it is important to note that the sum of forces \( F_S \) 
and the sum of force displacement \( F_{DS} \) is the basis for all the other variables.
It will be updated based on the inputs of users and the physics of the environment.
That again comes from the drone control script and the map configuration file.

\subsection*{Equations of motion:}
The equations of motion for the rigid body are given by:
\begin{align*}
    \ddot{P} &= \frac{F_S}{M} \\
    \tau &= F_S \times F_{DS} \\
    \dot{L} &= \tau_B \\
    \omega &= I^{-1} L \\
    \dot{q} &= \Omega(\omega)
\end{align*}
where:
\begin{tabbing}
	\hspace{1cm}\= \kill
	\> \( \Omega(\omega) \) is a axis quaternion \((||\omega||, |\omega|)\)
\end{tabbing}


\section{PID controller}

\subsection*{Additional Definitions:}
\begin{tabbing}
    \hspace{1cm}\= \kill
    \> \( P_d \) is the desired position in world frame. \\
    \> \( v_d \) is the desired velocity in world frame. \\
    \> \( q_d \) is the desired orientation. \\
    \> \( \omega_d \) is the desired angular velocity. \\
    \> \( e_P = P_d - P \) is the position error. \\
    \> \( e_v = v_d - v \) is the velocity error. \\
    \> \( e_q \) is the orientation error (quaternion difference). \\
    \> \( e_\omega = \omega_d - \omega \) is the angular velocity error. \\
    \> \( F_W \) is the total force in world frame. \\
    \> \( R(q, v) \) rotates the vector \( v \) by the quaternion \( q \). \\
    \> \( g \) is gravitational acceleration vector. \\
\end{tabbing}

\subsection*{PID Structure:}
Control is separated into translational and rotational components.

\subsection*{Translational Control (Position $\rightarrow$ Force):}

The desired total force in world frame is:
\begin{align*}
    F_W &= K_p^P e_P + K_d^P e_v + K_i^P \int e_P \, dt + M g
\end{align*}

This force is transformed into body frame:
\begin{align*}
    F_S &= R(q', F_W)
\end{align*}

\subsection*{Rotational Control (Orientation $\rightarrow$ Torque):}

The orientation error is computed using quaternion difference:
\begin{align*}
    e_q &= q_d \cdot q^{-1}
\end{align*}

The torque is defined by:
\begin{align*}
    \tau_B &= K_p^R \, \text{vec}(e_q) + K_d^R e_\omega + K_i^R \int \text{vec}(e_q) \, dt
\end{align*}

where \( \text{vec}(e_q) \) extracts the vector part of the quaternion.\\

This is just an example actual implementations may use more PID controllers 
or even just PI/PD or other combinations.

\subsection*{Force and Torque Allocation:}

The total force and torque are given by:
\begin{align*}
    F_S &= \sum_{i=1}^{n} F_i \\
    \tau &= \sum_{i=1}^{n} D_i \times F_i
\end{align*}
where:
\begin{tabbing}
    \hspace{1cm}\= \kill
    \> \( F_i \) is the force produced by the \( i \)-th motor. \\
    \> \( D_i \) is the position of the \( i \)-th motor relative to the center of mass. \\
\end{tabbing}

This can be written as a linear system:
\begin{align*}
    \begin{bmatrix}
        F_S \\
        \tau
    \end{bmatrix}
    =
    \begin{bmatrix}
        \mathds{1}_3 & \mathds{1}_3 & \cdots & \mathds{1}_3 \\
        [D_1]_\times & [D_2]_\times & \cdots & [D_n]_\times
    \end{bmatrix}
    \begin{bmatrix}
        F_1 \\
        F_2 \\
        \vdots \\
        F_n
    \end{bmatrix}
\end{align*}

\subsection*{Final Control Pipeline:}

At each timestep:
\begin{enumerate}
    \item Compute position and velocity error.
    \item Compute desired force \( F_S \).
    \item Compute orientation and angular velocity error.
    \item Compute torque \( \tau_B \).
    \item Solve for motor forces \( F_i \).
    \item Apply forces to update \( F_S \) and \( F_{DS} \).
\end{enumerate}

    \section{References}
    This math is based on the more in-depth document in the rocket simulation 
    project, available at the following link:
	\url{https://github.com/Lokestrom/rocketSim-code/blob/main/math.pdf}
\end{document}